% ============================================================
%  CHAPTER 8 — Conclusion
% ============================================================

\chapter{Conclusion}
\label{ch:conclusion}

% ============================================================
\section{Summary}
% ============================================================

This project delivered a complete, end-to-end hospital room and patient monitoring
system, structured around four tightly integrated subsystems: an ESP32\cite{esp32_datasheet} embedded
acquisition node, a Node.js/tRPC\cite{trpc} backend server, a React\cite{react} web dashboard, and a
cloud-hosted AI prediction layer.

Through rigorous quantitative evaluation, the system successfully demonstrated its capability to meet near real-time clinical monitoring constraints. The embedded firmware sustained an average sensor acquisition latency of 25\,ms, while the total edge-to-glass telemetry delay averaged $\approx$ 271\,ms over cloud infrastructure. The backend showed that a serverless, type-safe API can enforce clinical security requirements (RBAC, JWT verification, schema validation) with negligible added latency. The dashboard maintained a responsive 60\,FPS interface by intelligently throttling the heavy in-browser \texttt{face-api.js}\cite{faceapi} face recognition inference to a stable 1.6\,FPS loop. The AI layer demonstrated that simple statistical methods (Z-score anomaly detection and OLS linear regression) can deliver clinically meaningful early warnings with high sensitivity ($>90\%$), though deep learning approaches remain recommended for complex non-linear transitions.

The experimental results validate the architecture's core premise: a lightweight, decoupled system utilizing modern web technologies (WebSockets, Vercel\cite{vercel} edge deployment, Firebase\cite{firebase_rtdb}) can effectively serve as a highly available, robust clinical monitoring solution.

% ============================================================
\section{Achievement Against Objectives}
% ============================================================

All four functional objectives defined in Chapter~1 are addressed by the current
implementation. Continuous acquisition of temperature, humidity, heart rate, and
SpO\textsubscript{2} is designed for up to 1\,Hz transmission, constrained by sensor
limits. Data transmission to Firebase\cite{firebase_rtdb} RTDB operates on a 1\,s push interval, with
end-to-end latency dependent on network RTT. The web dashboard provides real-time
visualisation, alert management, and historical browsing. The AI layer classifies
clinical status and provides short-horizon vital sign forecasts based on recent
telemetry windows.

The non-functional objectives are addressed by design and verified by experiment: the system recovers
autonomously from Wi-Fi loss within 15 seconds (availability); the
NaN guard and IR threshold check prevent invalid records from reaching the
database (reliability); Firebase token verification and TLS-encrypted transport
protect clinical data (security); and the delta-threshold upload filter keeps
end-to-end latency dominated by network RTT rather than upload frequency
(performance), successfully achieving an average latency of $<300$\,ms.

% ============================================================
\section{Limitations}
% ============================================================

The most operationally significant limitation is the absence of a local offline
buffer in the embedded firmware: records generated during network outages are
permanently discarded, which is unacceptable in a clinical setting where data
continuity is a regulatory requirement. This should be the first priority in
any production hardening effort.

At the system level, the AI forecasting model's reliance on linear regression
constrains its utility to stable, slowly evolving physiological states. Rapid
non-linear transitions — such as sudden arrhythmias or acute respiratory events
— are not well captured, reducing the clinical value of the early-warning signal
in precisely the situations where it is most needed. This architectural limitation
motivates the migration to sequence-aware deep learning models as future work.

% ============================================================
\section{Future Work}
% ============================================================

The following directions are identified for future development of the complete system:

\begin{enumerate}
	\item \textbf{Local offline buffer with retry} (embedded): a LittleFS-backed
	      FIFO queue to preserve records during network outages and replay them on
	      reconnection, eliminating data loss.

	\item \textbf{NVS credential encryption} (embedded): enable the ESP32 native
	      NVS key partition to protect Firebase credentials at rest.

	\item \textbf{LSTM-based vital sign forecasting} (AI): replace the linear
	      regression extrapolator with a Long Short-Term Memory network trained on
	      real clinical time-series to capture non-linear physiological dynamics and
	      improve early-warning sensitivity.

	\item \textbf{Multi-room scalability} (backend + embedded): extend the data
	      model and tRPC API to support concurrent streams from multiple ESP32 nodes
	      deployed across different hospital rooms.

	\item \textbf{PWA with offline clinical records} (web): implement Progressive
	      Web App features for offline data caching and native push notifications,
	      reducing dependency on a persistent browser connection.

	\item \textbf{Motion artefact rejection} (embedded): add an accelerometer to
	      suppress SpO\textsubscript{2} and BPM readings during patient movement,
	      reducing AI false-positive rates.
\end{enumerate}

% ============================================================
\section{Deployment and Access}
% ============================================================

The web dashboard is deployed on Vercel\cite{vercel} and accessible at:
\begin{itemize}
	\item \textbf{URL:} \url{https://hospital-monioring.vercel.app/}
	\item \textbf{Test account:} \texttt{firas@mail.com}
	\item \textbf{Password:} \texttt{Firas*1234}
\end{itemize}
